% BHU PYQ -- Auto-generated & error-corrected
\documentclass[11pt,a4paper]{article}

\usepackage[a4paper,top=2.5cm,bottom=2.5cm,left=2.8cm,right=2.8cm,headheight=14pt]{geometry}
\usepackage{amsmath,amssymb}
\usepackage[T1]{fontenc}
\usepackage[utf8]{inputenc}
\usepackage{lmodern}
\usepackage{microtype}
\usepackage{fancyhdr}
\usepackage{enumitem}
\usepackage{hyperref}
\usepackage{lastpage}
\usepackage{parskip}

\hypersetup{colorlinks=true,urlcolor=black,linkcolor=black,
  pdftitle={STA/STB-503: Programming with C -- BHU PYQ},pdfauthor={Banaras Hindu University}}

\pagestyle{fancy}\fancyhf{}
\renewcommand{\headrulewidth}{0.4pt}\renewcommand{\footrulewidth}{0.4pt}
\fancyhead[L]{\small   \textit{Banaras Hindu University}}
\fancyhead[R]{\small   \textit{STA/STB-503: Programming with C}}
\fancyfoot[C]{\small Page  \thepage\ of \pageref{LastPage}}


\newlist{parts}{enumerate}{2}
\setlist[parts,1]{label=(\alph*),leftmargin=*,itemsep=5pt,topsep=3pt}
\setlist[parts,2]{label=(\roman*),leftmargin=*,itemsep=3pt,topsep=2pt}

\newcommand{\pts}[1]{\hfill   \textit{[#1]}}


\begin{document}


\begin{center}
  {\large\bfseries BANARAS HINDU UNIVERSITY}\\[2pt]
  {
\normalsize B.A./B.Sc. (Hons.) Semester V Examination, 2022-23}\\[6pt]
  \rule{0.8\linewidth}{0.5pt}\\[6pt]
  {\large\bfseries Paper No. STA/STB-503: Programming with C}\\[4pt]
  {
\normalsize Time: 3 Hours\hfill Maximum Marks: 70}
\end{center}
\rule{\linewidth}{0.4pt}
\smallskip



\noindent   \textit{Instructions: Answer any   \textbf{five} questions. All questions carry equal marks. Symbols have their usual meanings.}
\bigskip

\medskip


\noindent   \textbf{Question 1.}

\begin{parts}
  \item Explain type conversions in C expressions with examples. Discuss how operator precedence and associativity rules are applied. \pts{7}
  \item Determine the value of each of the following logical expressions, given   \texttt{int a = 5, b = 10, c = -1}:
  
\begin{parts}
    \item   \texttt{a > b \&\& a < c}
    \item   \texttt{a < b \&\& a > c}
    \item   \texttt{a == c || b > a}
    \item   \texttt{b > 15 \&\& c < 0 || a > 0}
  \end{parts}\pts{7}
\end{parts}

\medskip\rule{\linewidth}{0.2pt}

\medskip


\noindent   \textbf{Question 2.}

\begin{parts}
  \item Describe the different forms of the   \texttt{if-else} statement. How do they differ? Write a C program to calculate the mean of 10 numbers entered by the user. \pts{7}
  \item Write the output of the following program:

\begin{verbatim}
#include <stdio.h>
int main() {
    int x = 100;
    printf("%d
", 10 + x++ - 4);
    printf("%d
", 10 + --x + x);
    return 0;
}
\end{verbatim}\pts{7}
\end{parts}

\medskip\rule{\linewidth}{0.2pt}

\medskip


\noindent   \textbf{Question 3.}

\begin{parts}
  \item Define one-dimensional and two-dimensional arrays in C. Explain how a one-dimensional array is declared and initialized. Write a C program to find the transpose of a matrix. \pts{7}
  \item What is the output of the following program?

\begin{verbatim}
#include <stdio.h>
int main() {
    int m[] = {1, 2, 3, 4, 5};
    int x, y = 0;
    for (x = 0; x < 5; x++)
        y = y + m[x];
    printf("%d", y);
    return 0;
}
\end{verbatim}\pts{7}
\end{parts}

\medskip\rule{\linewidth}{0.2pt}

\medskip


\noindent   \textbf{Question 4.}

\begin{parts}
  \item Discuss how string variables are declared and initialized in C. Explain how strings are read from the terminal and written to the screen. Write a C program (without using library string functions) to copy one string into another and count the number of characters copied. \pts{7}
  \item Given   \texttt{char string1[10] = "abc"},   \texttt{string2[20] = "def"},   \texttt{string3[30]},   \texttt{string4[40]}, write the output of the following statements:

\begin{verbatim}
printf("%s", strcpy(string3, string1));
printf("%s", strcat(strcat(strcpy(string4, string1), "or"), string2));
printf("%d %d", strlen(string2) + strlen(string3), strlen(string4));
\end{verbatim}\pts{7}
\end{parts}

\medskip\rule{\linewidth}{0.2pt}

\medskip


\noindent   \textbf{Question 5.}

\begin{parts}
  \item Explain the different categories of functions in C with proper examples. Discuss the relevance of storage classes on the scope, visibility, and lifetime of variables. Write a C program that uses a function to sort an array of integers. \pts{7}
  \item What are user-defined functions? State three advantages of using functions. What is the purpose of a   \texttt{return} statement? \pts{7}
\end{parts}

\medskip\rule{\linewidth}{0.2pt}

\medskip


\noindent   \textbf{Question 6.}

\begin{parts}
  \item The Fibonacci numbers are defined recursively as: $F_1 = 1$, $F_2 = 2$, $F_n = F_{n-1} + F_{n-2}$ for $n > 2$. Write a C program containing a recursive function to generate and print the first 20 Fibonacci numbers. \pts{7}
  \item What is printed when the following program is executed?

\begin{verbatim}
#include <stdio.h>
int f(int);
int main() {
    printf("%d
", f(7));
    return 0;
}
int f(int x) {
    if (x <= 4) return x;
    return f(x - 1) + f(x - 2);
}
\end{verbatim}\pts{7}
\end{parts}

\medskip\rule{\linewidth}{0.2pt}

\medskip


\noindent   \textbf{Question 7.} With examples and programs, explain the following concepts in C. Give a typical use case for each:\pts{14}

\begin{parts}
  \item Structures
  \item Array of structures
  \item Nested structures
  \item Unions (include similarities and differences with structures, and rules for initializing)
\end{parts}

\vfill
\rule{\linewidth}{0.4pt}

\begin{center}\small
  \href{https://bhu-exam.vercel.app/}{Click here to visit our website}\\[4pt]
  \href{https://me-aryan.vercel.app/}{Click here to visit our contributor}\\[4pt]
  \href{https://quashan-me.vercel.app/}{Click here to visit our contributor}
\end{center}

\end{document}